\section{State of the Art in Industrial Agent Grounding}
\label{sec:agents}

Published industrial agents differ less in the language model they use than in where their
engineering evidence comes from and who holds authority over actions. We therefore classify
systems into five grounding categories (Fig.~\ref{fig:categories}) and code each included
system with AUF on what it actually uses at runtime (Table~\ref{tab:agent-systems}).

\begin{figure}[t]
\centering\footnotesize
\setlength{\tabcolsep}{2pt}
\begin{tabular}{@{}cl@{}}
\toprule
& Pipeline \\
\midrule
A & LLM $\rightarrow$ documents, manuals $\rightarrow$ recommendation \\
B & LLM $\rightarrow$ typed tools $\rightarrow$ industrial API \\
C & LLM $\rightarrow$ knowledge-graph retrieval $\rightarrow$ engineering context \\
D & LLM $\rightarrow$ proposal $\rightarrow$ engineering validation $\rightarrow$ execution \\
E & LLM $\rightarrow$ validation $\rightarrow$ dispatch gate $\rightarrow$ controller checks $\rightarrow$ process \\
\bottomrule
\end{tabular}
\caption{Grounding categories used to code industrial agent systems. Categories are cumulative
in authority, not in capability: a C system may lack typed tools.}
\label{fig:categories}
\end{figure}

\paragraph{A: Document-grounded.} Retrieval over manuals, procedures and logs
\cite{lewis2020rag} grounds explanations and troubleshooting. Identity (U1) and units (U3)
remain implicit in text, so answers can be fluent yet unverifiable, and actions have no
checkable target.

\paragraph{B: Typed-tool grounded.} The model calls declared functions on \opcua{}, AAS or
vendor APIs \cite{schick2023toolformer,yao2023react,patil2024gorilla,mcp}. Schemas provide
types and access (U2) and can constrain identifiers, but units, state legality and guards
are only as explicit as the exposed interface; tool transport does not establish equipment
truth.

\paragraph{C: Knowledge-graph grounded.} Retrieval over engineering graphs supplies identity,
relations and query pushdown (U1, U4, U8) \cite{edge2024graphrag,sarmah2024hybridrag,pan2024unifying,todkar2025semanticitot}.
Graphs built from engineering systems carry facts that text-derived graphs cannot, but
retrieval quality and the completeness of the mapped constraints become new dependencies.

\paragraph{D: Semantically validated.} Proposals are lifted into a checkable form and
validated against constraints and state models (U5, U6) before execution
\cite{knublauch2017shacl}. Validation converts a claim into a proposition with a truth
value, but only relative to the constraints that were modelled.

\paragraph{E: Controller-authoritative.} Validation is repeated at dispatch with fresh
evidence (U7), and the controller retains independent, current-state checks and final
actuator authority. This category separates supervisory reasoning from deterministic
control and is the only one compatible with functional-safety practice; it is also the
least reported with physical evaluation.

\begin{table*}[t]
\centering\footnotesize
\setlength{\tabcolsep}{3pt}
\caption{AUF coding of industrial AI-agent systems (0 absent, 1 prose/implicit, 2 explicit, 3 checked at runtime). Cat.: grounding category A--E (Fig.~\ref{fig:categories}).}
\label{tab:agent-systems}
\begin{tabular}{@{}llcccccccccccc@{}}
\toprule
System & Cat. & U1 & U2 & U3 & U4 & U5 & U6 & U7 & U8 & U9 & U10 & Valid. & Eval. \\
\midrule
\multicolumn{14}{c}{\textit{Full-text coding of included systems pending (Sec.~\ref{sec:method}).}} \\
\bottomrule
\end{tabular}
\end{table*}


The coding records, per system, which AUF capabilities are explicit or runtime-checked,
whether validation and controller authority exist, and the evaluation level. The
analysis questions are fixed in advance: which capabilities are most often absent (RQ4),
which category combinations co-occur with runtime validation (RQ5), and which studies
report measured effects of engineering semantics on grounding, action correctness or
context cost rather than demonstrations (RQ6). The number of coded systems
(\AgentSystemsCoded{} at the time of writing) and their distribution over categories
are regenerated from the released coding sheet.
